\documentclass[11pt,a4paper]{article}
\usepackage[margin=24mm]{geometry}
\usepackage{amsmath,amssymb,booktabs,graphicx}
\usepackage[hidelinks]{hyperref}
\usepackage{microtype}
\setlength{\parskip}{2pt}
\newcommand{\sig}[1]{\sigma_{#1}}
\newcommand{\code}[1]{\texttt{#1}}

\title{\textbf{The amplitude model as a generator}\\[2mm]
\large What \code{amp2609} is, where it may be used, and what it inherits}
\author{V.\ Kubarovsky\\ \small Jefferson Lab}
\date{18 September 2026}

\begin{document}
\maketitle

\begin{abstract}
\noindent
\code{exclurad\_py} now generates $\pi^0$ and $\eta$ electroproduction from a
model fitted directly to the world data at the level of the helicity amplitudes.
This note covers the join between the two: which parameter vector is installed,
what cross section is sampled, how far the sampling box reaches beyond the
region the model was fitted in, whether the sampled cross section stays
positive, and in what sense the radiative-correction loop is and is not closed.
The physics of the model itself is not repeated here; it is in
\code{\textasciitilde/rc\_iter1/tex/report.pdf}, and the radiative corrections
are in \code{doc/note/exclurad\_py\_note.pdf}.  This note is only the seam.
\end{abstract}

\section{What is installed}

\begin{center}
\begin{tabular}{ll}
\toprule
registered name & \code{pi0.amp2609}, \code{eta.amp2609}, \code{pi0n.amp2609}\\
parameter file & \code{exclurad\_py/models/amp2609\_par.npy}, 37 long\\
md5 & \code{abcb2f6ce7c87864386d900d82ad13ef}\\
origin & \code{\textasciitilde/rc\_iter1/runs/C1\_with\_compass/fitpar.npy}\\
fitted to & \code{pi0\_eta\_database.xlsx}, md5 \code{8d5b6a04206de590a186b0f205eea44c}\\
quality & $\chi^2/\mathrm{ndf} = 1.4018 = 2140.5/1527$, eleven sets, 1555 points\\
computed by & \code{\_amplitude\_fit2.py} (the 37-parameter layout)\\
\bottomrule
\end{tabular}
\end{center}

\noindent
The installed vector was verified byte-for-byte against the fit it claims to
come from.  Three normalisations were fitted alongside it --- 1.0586 for De
Masi, 0.9776 for the Hall A neutron, 1.0583 for Hall A 2016 --- and they are
\emph{not} applied here: they describe those experiments' normalisations, not
the physics.

\paragraph{\code{amp2026c} is withdrawn.}  It holds the previous fit,
\code{C0\_with\_compass}, which was made with a wrong kinematics for Hall A
2021: table I of PRL \textbf{127} 152301 heads its columns $0.36/0.48/0.60$ and
gives the per-setting $\langle x_B\rangle$ underneath, and the workbook had
taken the heading, putting four of the nine settings at the wrong $x_B$ by up to
$0.03$.  Anything generated with \code{amp2026c} after 2026-09-05 21:01 carries
that error.  The practical size of it is small --- between the two fits
$\chi^2/\mathrm{ndf}$ moves $1.3918 \to 1.4018$, $n_Q(H_T)$ $1.227 \to 1.260$,
and cross sections shift by a few percent --- so the events are not nonsense,
but they were made by a model that no longer exists and should be labelled, not
quietly kept.  The name must not be reused: an installed name is the only link
between a file on \code{wd\_14tb} and the physics that made it.

\section{What is sampled}

The model supplies nine structure functions in nb/GeV$^2$.  With an unpolarised target the generator samples
samples
\begin{equation}
\frac{d^2\sigma}{dt\,d\phi} = \frac{1}{2\pi}\Big[
\sig{T} + \epsilon\,\sig{L}
+ \epsilon\cos 2\phi\; \sig{TT}
+ \sqrt{2\epsilon(1+\epsilon)}\cos\phi\; \sig{LT}
+ h\sqrt{2\epsilon(1-\epsilon)}\sin\phi\; \sig{LT'}\Big],
\label{eq:phi}
\end{equation}
with $\phi$ in the Trento convention and $h$ the beam helicity.  This is the
decomposition the data were published in, so no conversion happens between the
fit and the generator.

Since 2026-09-18 the generator also samples the four structure functions a
\emph{longitudinally polarised target} needs.  With target polarisation
$\Lambda$, set by \code{ptarg} in the card, the bracket gains
\begin{equation}
+\ \Lambda\Big(\sqrt{2\epsilon(1+\epsilon)}\sin\phi\; \sig{LT}^{UL}
             + \epsilon\sin 2\phi\; \sig{TT}^{UL}\Big)
+\ h\Lambda\Big(\sqrt{1-\epsilon^2}\; \sig{T'}
              + \sqrt{2\epsilon(1-\epsilon)}\cos\phi\; \sig{LT'}^{LL}\Big),
\end{equation}
the first line being $A_{UL}$ and the second $A_{LL}$.  All four are fitted:
the eg1-dvcs moments of each are in the fit the parameters come from.

Three things about that are worth knowing.

\emph{The coefficients were measured, not assumed.}  Each candidate tensor
structure was contracted with the lepton tensor over a grid in $\phi$ and its
harmonic read off; the normalisation then follows from the same closure that
fixes \code{NU\_LT} and \code{NU\_LTP} in \code{core/hadronic.py}, and that
closure returns $-1$ for both of those, which is what makes the other four
trustworthy.  Born closure with $\Lambda$ is flat to $10^{-8}$ over both
helicities, $\Lambda \in \{0, \pm 1, 0.8\}$, nine $\phi$ and three kinematics.

\emph{The radiative correction is polarisation-dependent, correctly.}
\code{sigma\_F} contracts the hadronic tensor at the \emph{shifted} kinematics
rather than scaling a scalar Born, so the polarised terms are integrated over
photon emission like everything else, and the rotation of the photon frame
against the fixed target spin falls out of the basis vectors being built from
the actual momenta.  $\eta$ runs from 1.058 to 1.084 across $(h, \Lambda)$
where before there was one number.

\emph{The spin axis is the virtual photon, not the beam.}  The two differ by
the photon scattering angle, and the difference mixes a transverse target
asymmetry into $A_{UL}$.  The amplitude fit makes the same approximation, so
generator and model agree and are approximate in the same way; the angle is
the obvious next refinement, and it would require refitting eg1.

End to end, 8000 Born events at each of $\Lambda = \pm 0.8$ with the beam
helicity fixed give $\langle 2\sin\phi\rangle = +0.248 \pm 0.017$ against
$-0.275 \pm 0.018$.  The model gives $A_{UL}^{\sin\phi} = 0.30$ there, so the
modulation should be $0.8 \times 0.30 = 0.24$.

\emph{Nothing else changed.}  The polarised terms are gated on both a non-zero
$\Lambda$ and a model that supplies the keys, and \code{ptarg} defaults to
zero, so the vector-meson channels are bit-identical: born, $\eta$ and
$\sigma_{\rm obs}$ for \code{rho.diffrad\_vm}, \code{phi.diffrad\_vm},
\code{jpsi.diffrad\_vm} and \code{jpsi.diffrad} were compared against a
worktree at the previous commit over 48 values with zero differences.

\section{Where it may be used}

The model was fitted over
\[
Q^2 \in [1.15,\ 8.31]\ \mathrm{GeV}^2,\qquad
x_B \in [0.13,\ 0.60],\qquad
|t| \le 1.8\ \mathrm{GeV}^2 .
\]
The production card \code{cards/rad\_pi0\_2026.txt} samples $W^2 \in [4,12]$,
$Q^2 \in [1,5]$, $|t| \in [0.10,\ 2.0]$ at $E = 10.6$~GeV.  Scanning that box on
a grid of 62\,900 kinematic points:

\begin{center}
\begin{tabular}{lrr}
\toprule
 & points & share\\
\midrule
inside the fitted region & 50\,000 & 79.5\,\%\\
$|t| > 1.8$ & 8\,405 & 13.4\,\%\\
$x_B$ outside $[0.13, 0.60]$ & 3\,605 & 5.7\,\%\\
$Q^2 < 1.15$ & 3\,198 & 5.1\,\%\\
\bottomrule
\end{tabular}
\end{center}

\noindent
So about a fifth of the sampled box is extrapolation, and the largest part of it
is in $|t|$.  Two remarks on that.

The card stops at $Q^2 = 5$ while the fit now reaches $8.31$: the upper half of
the fitted range is not being used at all.  That is a card choice left from
before Hall A 2021 entered the fit, and it is worth revisiting --- that region is
now measured rather than extrapolated.

Extrapolation in $|t|$ is the least safe of the three, because the $t$
dependence is exactly what the form factors parametrise: eq.~(1) of the report
falls as $\exp[(b + b'\ln x_B)t]$ with no curvature, so beyond the data the
shape is an assumption, not a measurement.

\section{Positivity}

A generator cannot sample from a negative weight, and nothing in the fit
guarantees the bracket of eq.~(\ref{eq:phi}) stays positive: $\sig{TT}$,
$\sig{LT}$ and $\sig{LT'}$ enter with either sign and are not bounded by
$\sig{T}+\epsilon\sig{L}$ in the parametrisation.  It was therefore checked
rather than assumed.

Over the card's box, on 263\,808 points $(Q^2, W^2, |t|, \phi)$ with $\phi$ in
$7.5^\circ$ steps, the bracket never goes negative.  Its smallest value relative
to $\sig{T}+\epsilon\sig{L}$ is
\[
\min \frac{[\,\cdot\,]}{\sig{T}+\epsilon\sig{L}} = 0.141
\qquad\text{at } Q^2 = 1.00,\ x_B = 0.178,\ -t = 2.00,\ \phi = 180^\circ ,
\]
which is at the far corner of the extrapolated region, not inside the data.  The
margin is comfortable, but it is a property of this vector: it must be re-checked
whenever a new one is installed, and the scan above is cheap enough that there is
no excuse not to.

\section{The radiative-correction loop}

This is the part most easily got wrong, so it is worth stating exactly.

\paragraph{Inside the generator the loop is closed by construction.}
\code{core/rc.py:rc\_factor} takes a single \code{sfmodel} and passes the same
object to \code{born\_xsec} and to \code{sigma\_F}, the integral over real photon
emission.  There is no second model inside the radiative correction: whatever
the card names is used both for the Born cross section and for the soft and hard
photon integrals, and
\[
\eta = \sigma_{\rm obs}/\sigma_{\rm Born}
\]
is computed with that one model throughout.  Nothing needs fixing here.

\paragraph{On the data side the loop is at iteration one, and stopped there.}
The CLAS6 cross sections the model is fitted to are not the published ones: they
were re-corrected by
\[
\sigma_{\rm new} = \sigma_{\rm pub} / r, \qquad
r = \frac{\eta(\text{new model})}{\eta(\text{published model})},
\]
both $\eta$ computed with \code{exclurad\_py} so that the code cancels and only
the model difference survives (\code{compute\_rc.py}; published models
\code{pi0.vpk2013} with $v_{\rm cut} = 0.18$ and \code{eta.tableB1} with
$v_{\rm cut} = 0.094$).

The ``new model'' in that ratio is \textbf{\code{amp2021}}.  The model that now
describes those corrected data is \code{amp2609}, which is a different fit to a
larger and repaired database.  So the iteration that gives \code{rc\_iter1} its
name has been run once and not closed: the data carry the radiative correction
of an August model, and the model fitted to them is a September one.

The size of the inconsistency is not known and is directly computable: recompute
$r$ with \code{amp2609}, apply it, refit, and see how far the structure functions
and the parameters move.  Until that is done, absolute rates from the generator
are self-consistent \emph{within} the generator but rest on data corrected with a
different model.  This is the one open item in the chain, and it is a real one.

\section{What the generator inherits from the fit}

These are measured failures of the model, not caveats added for form.  Each is
documented with numbers in the report; here is what each means for generated
events.

\paragraph{$\sig{LT}$ is the model's worst-described quantity.}  It carries the
entire $\cos\phi$ modulation of eq.~(\ref{eq:phi}).  Left out of the fit, the
CLAS12 cross sections are predicted at 20.4 per point in $\sig{LT}$ against 2.39
in $\sig{U}$ and 3.51 in $\sig{TT}$; Hall A 2016 gives 4.7.  The $\phi$ shape of
the generated sample is therefore the least trustworthy thing it produces, and
any analysis that leans on the $\cos\phi$ moment should be aware of it.

\paragraph{A single power of $Q^2$ spans $1.2$ to $8.3$, and does not fit.}
Hall A 2016 and Hall A 2021 disagree with the model in opposite directions at the
two ends.  With Hall A 2021 in the fit the compromise costs the asymmetries.
Rates near the ends of the $Q^2$ range are less reliable than in the middle.

\paragraph{The neutron channel runs, and its normalisation is soft by a factor
two.}  \code{pi0n.amp2609} works and is a genuine prediction, but $n/p$ rests on
four points with 42 to 90\,\% errors, and a single constrained variant of the fit
--- tying the $t$-slopes of $\bar E_T^u$ and $\bar E_T^d$, which costs 194 units
of $\chi^2$ --- halves it, from $1.05$ to $0.52$, while the neutron data
themselves fit \emph{better}.  Treat neutron output as indicative.

\paragraph{Two parameters sit on a bound.}  The $t$-slopes of $H_T^u$ and of the
longitudinal block are both exactly zero, held there by the requirement that form
factors fall rather than grow with $|t|$.  The constraint is deliberate and
justified, but it means those slopes are limits, not measurements, and the
extrapolation in $|t|$ noted above starts from a boundary value.

\section{Reproducing a run}

A production run is defined entirely by its card.  For the record, the shape of
one:

\begin{center}
\begin{tabular}{ll}
\toprule
\code{pi0.amp2609} & model, and therefore the physics\\
\code{10.6} & beam energy, GeV\\
\code{200000} & events\\
\code{0} & \code{iborn}: 0 = full RC, 1 = Born only\\
\code{4.0 12.0} & $W^2$ range\\
\code{1.0 5.0} & $Q^2$ range\\
\code{-2.0 -0.10} & $t$ range\\
\code{0.50} & $v_{\rm cut}$\\
\code{0.80} & \code{ptarg}: longitudinal target polarisation, $-1\ldots+1$\\
\code{20260828} & seed\\
\code{/Volumes/wd\_14tb/mc/...} & output root\\
\bottomrule
\end{tabular}
\end{center}

\noindent
Run from \code{\textasciitilde/ampgen}, which holds the cards, the run scripts
and the logs, with

\begin{center}
\code{python3 -m exclurad\_py.run cards/<job>.txt --jobs 8}
\end{center}

\noindent
Two things about that directory are worth fixing.  It is not under version
control, so a card can change without leaving a trace, and the link between an
output directory on \code{wd\_14tb} and the card that produced it exists only by
naming convention.  For a generator whose whole value is reproducibility, that is
the weakest link in the chain --- weaker than anything in the physics above.

\section*{Where the rest is}

\begin{center}
\begin{tabular}{ll}
\toprule
the model and the fit & \code{\textasciitilde/rc\_iter1/tex/report.pdf}, 28 pages\\
the radiative corrections & \code{doc/note/exclurad\_py\_note.pdf}\\
the formulas & \code{docs\_src/rc\_formulas.pdf}\\
using the code & \code{doc/manual/exclurad\_py\_manual.pdf}\\
installing a new vector & \code{doc/HANDOVER\_amplitude\_model\_2026-09-05.md}\\
the fit, code and history & \url{https://github.com/vkubarovsky/rc_iter1}\\
\bottomrule
\end{tabular}
\end{center}

\end{document}
